% ============================================================
\chapter{Modul 11: NumPy -- Mathematik für Robotik}
% ============================================================

\section{Lektion 31: Warum NumPy?}

\begin{analogie}
Python-Listen können Zahlen speichern -- aber nicht damit rechnen. \texttt{[1,2,3] + [4,5,6]} ergibt \texttt{[1,2,3,4,5,6]}, nicht \texttt{[5,7,9]}. NumPy-Arrays machen das richtig -- und 100-mal schneller als Python-Schleifen.
\end{analogie}

Kinematik, Koordinatentransformationen, Filteralgorithmen, Bildverarbeitung -- alles in der Robotik basiert auf Matrixrechnung. NumPy ist das Werkzeug dafür.

\section{Lektion 32: Arrays erstellen und nutzen}

\begin{lstlisting}
import numpy as np

# Vektoren und Arrays
position  = np.array([1.5, 0.8, 0.0])         # x, y, theta
v_raeder  = np.array([0.30, 0.35])             # links, rechts
scan      = np.random.uniform(0.1, 3.0, 360)   # 360-Grad-Scan

# Rechnen: Operationen gelten fuer jedes Element
v_kmh = v_raeder * 3.6
print(f"Geschwindigkeiten in km/h: {v_kmh}")

# Scan auswerten
print(f"Naeherstes Hindernis: {scan.min():.2f} m")
print(f"Freieste Richtung:   {scan.max():.2f} m")
print(f"Mittlerer Abstand:   {scan.mean():.2f} m")

hindernisse = np.where(scan < 0.5)[0]   # Indices
print(f"Hindernisse in Richtungen: {hindernisse}")

# Matrizen: 2D-Arrays
rotation = np.array([
    [np.cos(np.pi/4), -np.sin(np.pi/4)],
    [np.sin(np.pi/4),  np.cos(np.pi/4)]
])
punkt = np.array([1.0, 0.0])
rotiert = rotation @ punkt   # Matrixmultiplikation
print(f"Rotierter Punkt: {rotiert.round(3)}")
\end{lstlisting}

\begin{merksatz}
\texttt{@} ist der Matrixmultiplikations-Operator in Python (seit 3.5). \texttt{*} ist elementweise Multiplikation. \texttt{np.dot(A, B)} tut dasselbe wie \texttt{A @ B}. In der Robotik ist \texttt{@} Standard.
\end{merksatz}

\section{Lektion 33: Koordinatentransformationen}

\begin{lstlisting}
def rotationsmatrix_2d(theta):
    """SO(2)-Rotationsmatrix fuer Winkel theta (Bogenmaß)."""
    c, s = np.cos(theta), np.sin(theta)
    return np.array([[c, -s], [s, c]])

def homogene_transformation(theta, tx, ty):
    """SE(2)-Transformationsmatrix (Rotation + Translation)."""
    c, s = np.cos(theta), np.sin(theta)
    return np.array([
        [c, -s, tx],
        [s,  c, ty],
        [0,  0,  1]
    ])

# Roboter bei (1.5, 0.8), Orientierung 45 Grad
theta = np.radians(45)
T_welt_roboter = homogene_transformation(theta, 1.5, 0.8)

# Sensor-Punkt im Roboter-Koordinatensystem
p_roboter = np.array([0.5, 0.0, 1.0])   # 0.5 m geradeaus

# Punkt in Weltkoordinaten transformieren
p_welt = T_welt_roboter @ p_roboter
print(f"Sensor-Punkt in Welt: ({p_welt[0]:.3f}, {p_welt[1]:.3f})")
\end{lstlisting}

\begin{praxis}
Homogene Koordinaten und SE(2)-Transformationen sind exakt die Mathematik aus dem KIT-Robotik-Curriculum (Kapitel 1). NumPy macht diese Berechnungen in Python direkt lösbar -- ein Array ist eine Matrix, \texttt{@} ist die Matrizenmultiplikation.
\end{praxis}

\begin{ros}
ROS2 nutzt intern die \texttt{tf2}-Bibliothek für Koordinatentransformationen -- dieselbe Mathematik wie oben, aber automatisiert für das gesamte Roboter-Koordinatensystem-Netz.
\end{ros}

\begin{uebung}[title=Projektaufgabe Modul 11: Odometrie-Berechnung]
Implementiere eine Odometrie-Simulation mit NumPy: Starte bei Position (0,0), Orientierung 0. Führe 100 Zeitschritte durch, in denen der Roboter mit $v = 0.3$ m/s und $\omega = 0.1$ rad/s fährt. Berechne mit homogenen Transformationen die Position nach jedem Schritt. Speichere alle Positionen in einem NumPy-Array.
\end{uebung}
